\hspace{3mm}

\centerline{\bf \Huge Домашнее задание.}
\centerline{\bf \Huge Сумма арифметической прогрессии}

\begin{flushright}
    \scriptsize{Даже самый быстрый глаз можно обмануть. \\Гото. Хантер}
\end{flushright}


\problem{1!} Найдите сумму 1 + 2 + 3 + … + 2000 + 2001 + 2002

\problem{2!} Найдите сумму 1 + 3 + 5 + … + 9995 + 9997 + 9999

\problem{3} На сколько сумма всех четных чисел от 2 до 10000 больше суммы всех нечетных чисел от 1 до 9999?

\problem{4} В первый день Шиншиллы решили 25 задач, а в последний - 975. Сколько дней Шиншиллы решали задачи, если в каждый последующий день они решали на одно и то же количество задач больше, чем в предыдущий, и общее количество решенных задач равно 10000?

\problem{5}\textbf{Бонусная задача.} 5 марта Шиншиллы выиграли 50 олимпиад, а 25 марта целых 10000 олимпиад. Сколько олимпиад выиграли Шиншиллы с 1 по 30 марта, если в каждый последующий день они выигрывали на одно и то же количество олимпиад больше, чем в предыдущий?